\documentclass[11pt]{article}
\usepackage{amsmath}
\usepackage{amssymb}
\usepackage{amsthm}
\usepackage{enumerate}
\usepackage[notcite,notref]{showkeys}
\usepackage{kotex}
\usepackage{graphicx}

\usepackage{url}
\usepackage{hyperref}

% PAGE SETTING ---------------------------------------

\usepackage[a4paper, margin=1in]{geometry}
\linespread{1.5}

% CUSTOM COMMAND -------------------------------------------

\renewcommand{\(}{\left(}
\renewcommand{\)}{\right)}
\renewcommand{\[}{\left[}
\renewcommand{\]}{\right]}

\newcommand{\inp}[2]{\langle #1,#2 \rangle}
\newcommand{\diag}{{\rm diag}}

\newcommand{\dist}{\text{dist }}
\newcommand{\supp}{\text{supp }}

\newcommand{\R}{\mathbb{R}}
\newcommand{\Rp}{\mathbb{R}_+}
\newcommand{\Rpp}{\mathbb{R}_{++}}
\newcommand{\C}{\mathbb{C}}
\newcommand{\N}{\mathbb{N}}
\newcommand{\Z}{\mathbb{Z}}
\newcommand{\Q}{\mathbb{Q}}

\newcommand{\ep}{\epsilon}
\newcommand{\pa}{\partial}

\newcommand{\mcC}{\mathcal{C}}
\newcommand{\mcH}{\mathcal{H}}
\newcommand{\mcT}{\mathcal{T}}
\newcommand{\mcV}{\mathcal{V}}
\newcommand{\mcG}{\mathcal{G}}
\newcommand{\mcE}{\mathcal{E}}
\newcommand{\mcW}{\mathcal{W}}

\newcommand{\suchthat}{ \; \big| \; }

%-----------------------------------------------------------------

\newcommand{\tags}[1]{#1}
\newenvironment{text-box}{}{}
\newcommand{\subheading}[1]{\noindent[ #1 ]}

% ----------------------------------------------------------------

\begin{document}
\title{This Is My LaTeX-to-Markdown Preset}
\date{2025-05-14}
\tags{}
\author{skwodnjs}

\maketitle

\section{Who am I?}

\subsection{Profile Image}

This is my profile image:

\begin{center}
    \includegraphics[width=0.3\textwidth]{../../../../assets/img/profile.jpg}
\end{center}

But it's not actually my cat.

\subsection{My Github Page}

My GitHub page is available at \url{https://skwodnjs.github.io/}. 
It contains mathematical knowledge I have studied and documented.
I first write in LaTeX, then convert it to Markdown.
I have written a Python script to automate this conversion.

\section{Github page?}

\subsection{How it works}

\begin{enumerate}
    \item Remove LaTeX commands related to metadata and document structure
    \item Convert \verb|\section| to \verb|#|, \verb|\subsection| to \verb|##|, and \verb|\subsubsection| to \verb|###|, \dots
    \item Wrap \verb|equation| and \verb|align| environments with \verb|$$|
    \item Replace \verb!|! and \verb!\|! inside \verb|$...$| with \verb|\vert| and \verb|\Vert| respectively
    \item Replace underscores \verb|_| inside \verb|$...$| with \verb|\_|
    \item Clean up \verb|\begin{text-box}| and add \verb|{: .text-box}|
    \item Convert \verb|\subheading{...}| to Markdown-style headings
    \item Process image-related commands such as \verb|\includegraphics| and convert to Markdown format
    \item Trim leading/trailing whitespace and reduce consecutive blank lines to a single blank line
\end{enumerate}

\subsection{Example}

The \verb|\subheading| command is used for emphasis.

\subheading{proof}

The \verb|text-box| environment can be used to group and highlight explanations.

\begin{text-box}
    A collection $\tau$ of subsets of $X$ is called a topology on $X$ if \\
    (i) $X, \emptyset \in \tau$. \\
    (ii) The union of any collection of sets in $\tau$ is also in $\tau$. \\
    (iii) The finite intersection of sets in $\tau$ is also in $\tau$.
\end{text-box}

\end{document}